\chapter*{СВЕДЕНИЯ}                         % Заголовок
\addcontentsline{toc}{chapter}{ОБЩИЕ СВЕДЕНИЯ}    % Добавляем его в оглавление

Производственная практика проводилась в период с «05» февраля 2024
по «19» мая 2024~г. 

Руководителем практики от профильной организации являлась ст.~преп. М.~С.~Анисимова.

Практика проходила на кафедре Инфокоммуникационных технологий НИТУ МИСИС.

Содержание индивидуального задания: исследование влияния эмпатии в больших языковых моделях при их использовании в процессе принятия решений в теории игр, разработка экспериментальной методологии для анализа влияния эмоций в играх с торгом, подбор проприетарных и свободных моделей искусственного интеллекта, проведение экспериментов в многоязычных средах с разными эмоциональными состояниями, анализ результатов.

Утвержденный календарный график проведения практики:

% Please add the following required packages to your document preamble:
% \usepackage{graphicx}
\begin{table}[ht!]
\begin{tabular}{|p{0.2\linewidth}|p{0.75\linewidth}|}
\hline
Сроки   работы & Наименование видов работ                                                                                                                                                  \\
\hline
05.02          & Прохождение   рабочего инструктажа. Получение индивидуального задания                                                                                                     \\
\hline
06.02–17.02    & Сбор и   анализ необходимой документации для проведения анализа предметной области исследования                                                                           \\
\hline
19.02–9.03     & Сбор и   анализ теоретических сведений. Изучения комплекса методик и инструментальных   средств, использование которых позволяет решать поставленные задачи по   практике \\
\hline
11.03–11.05    & Выполнение индивидуального   задания                                                                                                                                      \\
\hline
13.05–16.05    & Обоснование   полученных результатов и написание отчета                                                                                                                   \\
\hline
до 19.05       & Защита   отчета по практике    \\                                                                                                                                         
\hline
\end{tabular}%
\end{table}

\chapter*{ВВЕДЕНИЕ}                         % Заголовок
\addcontentsline{toc}{chapter}{ВВЕДЕНИЕ}    % Добавляем его в оглавление
Симуляция человеческого поведения является признанной практикой, 
как в социологии, так и в экономике{\cite{ziems2024can}. Однако, конвенциональные методы 
симуляции ограничены в моделировании всех существенных аспектов 
человеческих взаимодействий \cite{xie2024can}. Последние разработки в сфере искусственного 
интеллекта, в частности большие языковые модели (БЯМ, LLM), 
использовали инструменты агентного моделирования. \cite{openai2023gpt4,wu2023autogen, chen2023agentverse, park2023generative}. Большинство 
исследователей использовали LLM в различных игровых теоретических 
условиях, что является наиболее естественным способом исследования 
лежащих в основе экономических и социальных причин человеческого 
поведения. \cite{akata2023playing,horton2023large,aher2023using}.

Тем не менее, большинство существующих исследований исходят из 
недоказанной предпосылки, что агенты LLM ведут себя аналогично людям в 
моделировании. Более того, они игнорируют критический аспект 
эмоционального человеческого поведения, который существенно влияет на 
процесс принятия решений \cite{lerner2015emotion}. Чтобы устранить эти пробелы, была разработана 
новая методология для изучения соответствия между поведением LLM и 
человека на основе их сравнения в различных эмоциональных состояниях. В 
частности, мы сосредоточились на пяти базовых эмоциях: гнев, грусть, 
счастье, отвращение и страх, которые были выбраны на основе хорошо 
установленной классификации Пола Экмана ~\cite{ekman1999basic}. Мы предлагаем новый подход к 
наделению LLM эмоциями \cite{akata2023playing, brookins} с помощью эмоционального промптинга  \cite{li2023large}. 
исследуя его влияние на соответствие процессу принятия решений человека в 
играх с торгом.

Цель данной работы: исследование влияния эмпатии в больших 
языковых моделях при их использовании в процессе принятия решений в 
теории игр, разработка экспериментальной методологии для анализа влияния 
эмоций в играх с торгом, подбор проприетарных и свободных моделей 
искусственного интеллекта, проведение экспериментов в многоязычных 
средах с разными эмоциональными состояниями, анализ результатов.

%\newcommand{\actuality}{}
%\newcommand{\progress}{}
%\newcommand{\aim}{{\textbf\aimTXT}}
%\newcommand{\tasks}{\textbf{\tasksTXT}}
%\newcommand{\novelty}{\textbf{\noveltyTXT}}
%\newcommand{\influence}{\textbf{\influenceTXT}}
%\newcommand{\methods}{\textbf{\methodsTXT}}
%\newcommand{\defpositions}{\textbf{\defpositionsTXT}}
%\newcommand{\reliability}{\textbf{\reliabilityTXT}}
%\newcommand{\probation}{\textbf{\probationTXT}}
%\newcommand{\contribution}{\textbf{\contributionTXT}}
%\newcommand{\publications}{\textbf{\publicationsTXT}}

%
\textbf{Актуальность} связана с необходимостью выравнивания LLM в более сложных стратегических и этических средах, включая мульти-агентные системы.

% {\progress}
% Этот раздел должен быть отдельным структурным элементом по
% ГОСТ, но он, как правило, включается в описание актуальности
% темы. Нужен он отдельным структурынм элемементом или нет ---
% смотрите другие диссертации вашего совета, скорее всего не нужен.

{\aim} данной работы является исследование влияния эмпатии в больших
языковых моделях при их использовании в процессе принятия решений в
теории игр, разработка экспериментальной методологии для анализа влияния
эмоций в играх с торгом, подбор проприетарных и свободных моделей
искусственного интеллекта, проведение экспериментов в многоязычных
средах с разными эмоциональными состояниями, анализ результатов.

Для~достижения поставленной цели необходимо было решить следующие {\tasks}:
\begin{enumerate}[beginpenalty=10000] % https://tex.stackexchange.com/a/476052/104425
  \item Рассмотреть теоретические аспекты влияния эмпатии в контексте теории игр.
  \item Разработать фреймворк для проведения экспериментов для оценки влияния эмоций в больших языковых моделях на примере игр <<Диктатор>>, <<Дилемма узника>> и <<Общественное благо>>.
  \item Провести эксперименты, связанные с эмоциональным воздействием, в исследуемых играх.
  \item Выявить модели, наиболее соответствующие выравниванию с эмоциональным поведением человека.
\end{enumerate}


{\novelty}
%\begin{enumerate}[beginpenalty=10000] % https://tex.stackexchange.com/a/476052/104425
Была проведена оценка влияний эмпатии на этическое и эмоциональное поведение LLM в контексте теории игр.
%\end{enumerate}

{\influence} Практическая значимость эмпатии для выравнивания в крупных языковых моделях заключается в её возможности увеличить эффективность в реальных жизненных ситуациях, не имеющих одногозначного решения. 

{\methods} В работе применялись методы теории игр, статистический анализ данных, методы управления большими языковыми моделями, общенаучные методы анализ и синтез.

%{\defpositions}
%\begin{enumerate}[beginpenalty=10000] % https://tex.stackexchange.com/a/476052/104425
%  \item Первое положение
%  \item Второе положение
%  \item Третье положение
%  \item Четвертое положение
%\end{enumerate}
%В папке Documents можно ознакомиться с решением совета из Томского~ГУ
%(в~файле \verb+Def_positions.pdf+), где обоснованно даются рекомендации
%по~формулировкам защищаемых положений.

%{\reliability} полученных результатов обеспечивается \ldots \ Результаты находятся в соответствии с результатами, полученными другими авторами.
 % Характеристика работы по структуре во введении и в автореферате не отличается (ГОСТ Р 7.0.11, пункты 5.3.1 и 9.2.1), потому её загружаем из одного и того же внешнего файла, предварительно задав форму выделения некоторым параметрам

%\textbf{Объем и структура работы.} Диссертация состоит из~введения,
%\formbytotal{totalchapter}{глав}{ы}{}{},
%заключения и
%\formbytotal{totalappendix}{приложен}{ия}{ий}{}.
%% на случай ошибок оставляю исходный кусок на месте, закомментированным
%Полный объём диссертации составляет  \ref*{TotPages}~страницу
%с~\totalfigures{}~рисунками и~\totaltables{}~таблицами. Список литературы
%содержит \total{citenum}~наименований.
%
%Полный объём диссертации составляет
%\formbytotal{TotPages}{страниц}{у}{ы}{}, включая
%\formbytotal{totalcount@figure}{рисун}{ок}{ка}{ков} и
%\formbytotal{totalcount@table}{таблиц}{у}{ы}{}.
%Список литературы содержит
%\formbytotal{citenum}{наименован}{ие}{ия}{ий}.
